\section{Faithful finite diagram envelopes}
\label{sec:envelopes}

The finite idempotent sets $\mathcal P_2$ and $\mathcal P_3$ specify
only finitely many commutation and multiplication identities.  We now
realize all these identities in a finitely presented overgroup, without
collapsing the common full-identity copy of the input group.  A faithful
model supplies the injectivity; its own presentation need not be finite.

\begin{proposition}\label{prop:finite-envelope}
For every finitely presented group $U$ there are a finitely presented
group $\widehat W$ and homomorphisms
\[
 k_{\ell,E}:U\longrightarrow\widehat W
 \qquad(\ell\in\{2,3\},\ E\in\mathcal P_\ell)
\]
such that $k_{2,I}=k_{3,I}$ is injective.  For each $\ell$, whenever
$E,E',E+E'\in\mathcal P_\ell$ and $EE'=E'E=0$, the subgroups
$k_{\ell,E}(U)$ and $k_{\ell,E'}(U)$ commute elementwise, and
\[
 k_{\ell,E+E'}(u)=k_{\ell,E}(u)k_{\ell,E'}(u)
 \qquad(u\in U).
\]
\end{proposition}

\begin{proof}
Fix a finite presentation $U=\langle X\mid\mathcal R\rangle$.
For every pair $(\ell,E)$ take a copy of this presentation, with
generators $x_{\ell,E}$ for $x\in X$.  Define $\widehat W$ by these
copy presentations and the following additional relations:
\begin{align}
 x_{2,I}&=x_{3,I} &&(x\in X),\label{eq:envelope-common}\\
 [x_{\ell,E},y_{\ell,E'}]&=1 &&(x,y\in X),\label{eq:envelope-commute}\\
 x_{\ell,E+E'}&=x_{\ell,E}x_{\ell,E'} &&(x\in X).
 \label{eq:envelope-sum}
\end{align}
The last two families run over both values of $\ell$ and every
orthogonal pair $E,E'$ whose members and sum lie in
$\mathcal P_\ell$.  All three lists are finite.  The indices are actual
idempotents in $D_\ell$, not separately tagged copies from its three
matrix subalgebras.  Thus these relations include the orthogonal sums
between different subalgebras as well as those within one subalgebra.
In particular, the zero idempotent is included; its sum relation with
itself gives $x_{\ell,0}=x_{\ell,0}^2$, forcing
$x_{\ell,0}=1$.

The copy presentations give homomorphisms $k_{\ell,E}:U\to\widehat W$.
Equation~\eqref{eq:envelope-commute} implies that their corresponding
image subgroups commute elementwise.  It follows that
$u\mapsto k_{\ell,E}(u)k_{\ell,E'}(u)$ is a homomorphism.  By
Equation~\eqref{eq:envelope-sum}, it agrees with $k_{\ell,E+E'}$ on
$X$, and therefore on all of $U$.  Similarly,
Equation~\eqref{eq:envelope-common} gives $k_{2,I}=k_{3,I}$.

To prove that this common map is injective, we construct a faithful
model.  The same group-level enlargement is used in the
companion~\cite[Section~3.1]{OpenAIOvergroups}; we verify all of its needed algebraic
properties here.  For any group $V$
and $\ell\in\{2,3\}$, work in the unital algebra
\[
 A_\ell(V)=D_\ell\otimes_{\F_\ell}\F_\ell[V],
\]
where $\F_\ell[V]$ is the group algebra with basis indexed by $V$.
For $E\in\mathcal P_\ell$ define
\begin{equation}\label{eq:envelope-model}
 j_{\ell,E}(v)=(I-E)\otimes1+E\otimes v
 \qquad(v\in V).
\end{equation}
Since $E$ and $I-E$ are orthogonal idempotents, multiplication gives
\[
 j_{\ell,E}(v)j_{\ell,E}(w)=j_{\ell,E}(vw),
 \qquad j_{\ell,E}(1)=1.
\]
Thus these elements are units, with inverses $j_{\ell,E}(v^{-1})$,
and $j_{\ell,E}$ is a homomorphism.

For $E,E',E+E'\in\mathcal P_\ell$ with $EE'=E'E=0$, rewrite
Equation~\eqref{eq:envelope-model} as $1+E\otimes(v-1)$.  Thus
\[
 \bigl(1+E\otimes(v-1)\bigr)
 \bigl(1+E'\otimes(w-1)\bigr)
 =1+E\otimes(v-1)+E'\otimes(w-1),
\]
because $EE'=0$.  Reversing the order gives the same answer since
$E'E=0$.  Hence the images of $j_{\ell,E}$ and $j_{\ell,E'}$
commute elementwise, even if $v$ and $w$ do not commute.  Setting
$w=v$ gives
\[
 j_{\ell,E}(v)j_{\ell,E'}(v)=j_{\ell,E+E'}(v).
\]
Moreover, $j_{\ell,I}(v)=I\otimes v$ is injective: the decomposition
$A_\ell(V)=\bigoplus_{v\in V}D_\ell\otimes v$ distinguishes these
elements, since $I\ne0$.

Let $W_\ell(V)$ be the subgroup of $A_\ell(V)^\times$ generated by
all images $j_{\ell,E}(V)$.  First form $W_2(U)$, and then form
$W_3(W_2(U))$.  Write $j_E^{(2)}$ for the first-stage maps and
$j_E^{(3)}$ for the second-stage maps.  The group $W_3(W_2(U))$
receives the two families
\[
 \alpha_E=j_I^{(3)}\circ j_E^{(2)}\quad(E\in\mathcal P_2),
 \qquad
 \beta_E=j_E^{(3)}\circ j_I^{(2)}\quad(E\in\mathcal P_3).
\]
Postcomposition preserves the first family's identities, and
precomposition preserves the second family's identities.  Their
full-identity maps coincide with the injection
$j_I^{(3)}\circ j_I^{(2)}$.  Notice that the second group algebra is
formed from the abstract group $W_2(U)$; no map between algebras of
different characteristics is involved.

These families satisfy every defining relation of $\widehat W$.
They therefore induce a homomorphism
$\Phi:\widehat W\to W_3(W_2(U))$ for which
\[
 \Phi\circ k_{2,I}=j_I^{(3)}\circ j_I^{(2)}.
\]
The right-hand side is injective, so $k_{2,I}$ is injective as claimed.
\end{proof}

We can now bring the entire finite diagram back into a universal group.

\begin{corollary}\label{prop:self-map}
Let $U$ be a finitely presented group containing an isomorphic copy of
every finitely presented group.  There are an injective endomorphism
$f:U\to U$ and homomorphisms
$h_{\ell,E}:U\to U$ for $\ell\in\{2,3\}$ and
$E\in\mathcal P_\ell$, with $h_{2,I}=h_{3,I}=f$, satisfying all the
orthogonal commutation and sum identities of
Proposition~\ref{prop:finite-envelope}.
\end{corollary}

\begin{proof}
Apply Proposition~\ref{prop:finite-envelope} to $U$.  Since
$\widehat W$ is finitely presented, universality supplies an embedding
$\kappa:\widehat W\to U$.  Set
\[
 h_{\ell,E}=\kappa\circ k_{\ell,E},
 \qquad f=\kappa\circ k_{2,I}=\kappa\circ k_{3,I}.
\]
The map $f$ is injective, and postcomposition by $\kappa$ preserves
every required identity and commutation relation.
\end{proof}
